\section{Related work}
\label{sec:related}

Three literatures meet on this question and none of them owns it outright. We group them, say
what each established, and end with what is left for us.

\subsection{Labor displacement and the public finances}

The proposition that labor-displacing automation is a fiscal event is established, and this
paper does not claim it. \citet{PriceSuresh2026} measure it in a budget model, putting most
federal revenue on individual or payroll taxes and roughly two thirds directly on labor, and
varying reemployment and the pricing of AI services, which are the same two quantities this paper
calls the retained wage share and the effective rate on capital. Their exercise holds output
constant and carries no household balance sheets. \citet{IeongSaputraManiarCheng2026} give a
four-channel account of how such a transition erodes a tax base, and \citet{Barhoumi2026} states
the mechanism directly in a workshop synthesis. We treat these as connected work rather than as
competition.

\citet{KorinekLockwood2026} ask what an optimal tax system should do across the stages of such a
transition. That is the normative counterpart of our question and it is not the same question: we
measure what the existing system reaches, and the gap between what is optimal and what is
reachable is the space this paper works in. \citet{FalkTsoukalas2026} reach a stronger conclusion
first, that because each firm captures the full cost saving from automation while bearing only
part of the demand loss, capital income taxes cannot resolve the externality and a Pigouvian
automation tax can. Our contribution there is calibration rather than verdict: we do not argue
against their instrument, we measure how far the conventional one falls short and why.
\citet{CasasTorres2024} established the base-shift mechanism in general equilibrium before this
paper, and \citet{AcemogluManeraRestrepo2020} supplies the expensing result we take as derived.
\citet{Huckfeldt2022} gives the wage loss on reemployment behind our retained wage share, fitted
on the displacement vintages of \citet{BLSWorkerDisplacement}, and \citet{Brollo2024} supplies the
one measured policy effect we rely on.

\subsection{Effective rates on capital income, and what erodes them}

The components of an effective marginal rate are well established separately.
\citet{Auerbach2017} sets out the treatment of expensing and the cash-flow base.
\citet{CBO2014} is the published measured benchmark our assembly is checked against.
\citet{TorslovWierZucman2023} measure the share of profit booked in lower-taxed jurisdictions,
which is the erosion channel that has received by far the most attention. The share of the return
that is economic rent has two incompatible readings, \citet{Barkai2020} and
\citet{KarabarbounisNeiman2014}, which are two accounts of what happened to the corporate sector
rather than two estimates of one quantity, so we report both and never average them.

The strand nearest our finding is the ownership literature.
\citet{RosenthalAustin2016}, \citet{RosenthalBurke2020} and \citet{RosenthalMucciolo2024} have
documented the long decline in the share of US corporate stock held in taxable accounts, and
\citet{Gravelle2022} the deferral and step-up that erode what remains of the shareholder layer.
Their question is how much of the shareholder tax base survives. Ours is adjacent and, as far as
we have located, not asked: what that shrinking base implies for the feasibility of replacing a
wage tax with a capital tax, and whether the same ownership structure that erodes the base also
governs who absorbs the shift. We take their parameters as inputs and carry the question forward.

\subsection{Household balance sheets and incidence}

The household side rests on the relationship between income shocks, liquid buffers and distress
\citep{MianRaoSufi2013,GanongNoel2020,BhuttaBrickerDettlingKelliherLaufer2019,MerikullRoom2020},
with the stress-testing template from \citet{AlbaceteFessler2010}.
\citet{BartscherKuhnSchularickSteins2020} trace inequality, household debt and fragility
together over the post-war period, and \citet{Bayraktar2026} models automation, income incidence
and capital accumulation in incomplete markets, deriving a growth threshold from a model where we
measure an incidence boundary from survey microdata. Our exercise is neither a debt-service
stress test in their sense nor a model-based threshold: it is scenario arithmetic on measured
holdings, and we label it as such wherever it appears.

The concentration of equity ownership itself is measured in the distributional financial accounts
\citep{FRBDFA}, and the wealth effect we use to price the bust direction is from
\citet{ChodorowReichNenovSimsek2021}.

\subsection{What is ours}

\begin{quote}
The fiscal mechanism is established and the ownership parameters are published. What this paper
adds is the condition stated as a boundary rather than a comparison; both of its sides assembled
from sourced components and independently rebuilt; a ranking of the levers by measured influence
that ranks the ownership channel above the profit-shifting one, with the sensitivity of that
ranking reported; and the demonstration that one ownership fact
governs the revenue side, the household side and the transmission side at once.
\end{quote}
